%%%%%%%%%%%%%%%%%%%%%%%%%%%%%%%%%%%%%%%%%%%%%%%%%%%%%%%%%%%%%%%%%%%%%%%%%%%%%%%%%%
\begin{frame}[fragile]\frametitle{}
\begin{center}
{\Large Making of a Sandwich}
\end{center}
\end{frame}

%%%%%%%%%%%%%%%%%%%%%%%%%%%%%%%%%%%%%%%%%%%%%%%%%%%%%%%%%%%
% Every slide in this section applies ONE prompt technique to the same task
% so the audience can isolate exactly what each technique contributes.
%%%%%%%%%%%%%%%%%%%%%%%%%%%%%%%%%%%%%%%%%%%%%%%%%%%%%%%%%%%

\begin{frame}[fragile]\frametitle{Basic Prompt}

\textbf{Prompt:} \emph{Explain how to make a peanut butter and jelly sandwich.}

\medskip
\textbf{What happens:} The AI gives a generic answer, it will assume standard ingredients and a general audience. This is fine for casual use, but the output can be inconsistent if you run it multiple times.

\medskip
\textbf{Lesson:} Even a simple prompt works, but gives you \emph{no control} over the style, audience, or constraints of the answer.

\end{frame}

%%%%%%%%%%%%%%%%%%%%%%%%%%%%%%%%%%%%%%%%%%%%%%%%%%%%%%%%%%%
\begin{frame}[fragile]\frametitle{Adding Roles (Persona Prompting)}

\textbf{Prompt:} \emph{As a chef, explain to your assistant how to make a peanut butter and jelly sandwich.}

\medskip
\textbf{What changes:} The AI now adopts the voice and vocabulary of a professional chef. The tone becomes instructional and precise rather than casual.

\medskip
\textbf{Lesson:} Giving the AI a \textbf{role} (``Act as a chef / lawyer / teacher'') immediately shapes the tone, vocabulary, and depth of the response, without changing the core question.

\end{frame}

%%%%%%%%%%%%%%%%%%%%%%%%%%%%%%%%%%%%%%%%%%%%%%%%%%%%%%%%%%%
\begin{frame}[fragile]\frametitle{Adding Constraints}

\textbf{Prompt:} \emph{Make a nut-free version of the sandwich due to a customer's nut allergy.}

\medskip
\textbf{What changes:} The AI now replaces peanut butter with an alternative (e.g., sunflower seed butter or jam-only) and explicitly avoids any nut-based ingredient.

\medskip
\textbf{Lesson:} \textbf{Constraints} (``avoid X'', ``must include Y'', ``limit to N words'') are powerful control levers. The AI will respect them if you state them clearly. Vague constraints lead to vague compliance.

\end{frame}


%%%%%%%%%%%%%%%%%%%%%%%%%%%%%%%%%%%%%%%%%%%%%%%%%%%%%%%%%%%
\begin{frame}[fragile]\frametitle{Adding Examples (Few-Shot)}

\textbf{Prompt:} \emph{Create two unique variations of the classic sandwich.
Example 1, Banana Nut Crunch: [description].
Example 2, Triple Berry Blast: [description]. Now create a third variation.}

\medskip
\textbf{What changes:} By showing two examples first, you teach the AI the \textbf{format, style, and level of detail} you expect, without writing a long list of instructions.

\medskip
\textbf{Lesson:} This is called \textbf{Few-Shot Prompting}, giving the AI a few examples to pattern-match from. It is one of the most reliable ways to get consistently formatted outputs.

\end{frame}

%%%%%%%%%%%%%%%%%%%%%%%%%%%%%%%%%%%%%%%%%%%%%%%%%%%%%%%%%%%
\begin{frame}[fragile]\frametitle{Adding Contextual Information}

\textbf{Prompt:} \emph{As the head chef at ``The Sandwich Haven'', guide your new assistant to create specials for the menu.}

\medskip
\textbf{What changes:} The AI now understands the \emph{setting} (a real restaurant, not a home kitchen), the \emph{relationship} (head chef to assistant), and the \emph{purpose} (menu specials, not everyday cooking).

\medskip
\textbf{Lesson:} \textbf{Context} narrows the creative space for the AI. The richer the context, the more targeted and useful the output becomes.

\end{frame}

%%%%%%%%%%%%%%%%%%%%%%%%%%%%%%%%%%%%%%%%%%%%%%%%%%%%%%%%%%%
\begin{frame}[fragile]\frametitle{Incorporating Feedback (Iterative Prompting)}

\textbf{Prompt:} \emph{Improve the sandwich based on customer feedback: customers want less sweetness and something more creative.}

\medskip
\textbf{What changes:} The AI refines its previous output based on specific, actionable feedback rather than starting from scratch.

\medskip
\textbf{Lesson:} Prompt Engineering is an \textbf{iterative process}, you rarely get the perfect answer on the first try. Treating your conversation with the AI like a feedback loop gives progressively better results.

\end{frame}

%%%%%%%%%%%%%%%%%%%%%%%%%%%%%%%%%%%%%%%%%%%%%%%%%%%%%%%%%%%
\begin{frame}[fragile]\frametitle{Time Constraints and Prioritization}

\textbf{Prompt:} \emph{Prepare an alternative fruit-based version for testing within a tight deadline, focus on the most important steps only.}

\medskip
\textbf{What changes:} The AI prioritises and condenses, it drops elaboration and gives you a quick, actionable plan.

\medskip
\textbf{Lesson:} You can encode \textbf{business constraints} directly into your prompt: urgency, scope, priority. The AI will adapt accordingly.

\end{frame}

%%%%%%%%%%%%%%%%%%%%%%%%%%%%%%%%%%%%%%%%%%%%%%%%%%%%%%%%%%%
\begin{frame}[fragile]\frametitle{Incorporating Multidisciplinary Knowledge}

\textbf{Prompt:} \emph{Use food presentation and garnishing techniques to describe how to make a visually appealing sandwich for a food photography session.}

\medskip
\textbf{What changes:} The AI draws on knowledge from \emph{multiple domains}, culinary arts, visual presentation, and photography staging.

\medskip
\textbf{Lesson:} LLMs hold knowledge from many fields simultaneously. A good prompt can \textbf{combine domains} in ways that would take a human significant research time to replicate.

\end{frame}

%%%%%%%%%%%%%%%%%%%%%%%%%%%%%%%%%%%%%%%%%%%%%%%%%%%%%%%%%%%
\begin{frame}[fragile]\frametitle{Addressing Dietary Preferences}

\textbf{Prompt:} \emph{Prepare a fully vegan version using only plant-based alternatives for all ingredients.}

\medskip
\textbf{What changes:} The AI substitutes every non-vegan ingredient (honey, dairy butter) with appropriate alternatives and justifies each substitution.

\medskip
\textbf{Lesson:} Constraints based on \textbf{values or requirements} (vegan, halal, gluten-free, accessible, budget-friendly) are perfectly understood by modern LLMs, you just need to state them.

\end{frame}

%%%%%%%%%%%%%%%%%%%%%%%%%%%%%%%%%%%%%%%%%%%%%%%%%%%%%%%%%%%
\begin{frame}[fragile]\frametitle{Self-Criticism (Reflexion / Self-Refinement)}

\textbf{Prompt:} \emph{Explain how to make a peanut butter and jelly sandwich. Please re-read your above response. Are there any mistakes or omissions? If so, identify and correct them.}

\medskip
\textbf{What changes:} The AI reviews its own output for logical errors, missing steps, or unclear instructions, and corrects them.

\medskip
\textbf{Lesson:} You can instruct the AI to \textbf{self-critique} its answers. This technique, sometimes called \emph{Reflexion}, produces significantly more accurate and complete outputs, especially for complex tasks.

\note{Self-refinement is particularly valuable for educators writing rubrics or assessment criteria: ask the AI to check its own draft for gaps before you review it.}

\end{frame}

%%%%%%%%%%%%%%%%%%%%%%%%%%%%%%%%%%%%%%%%%%%%%%%%%%%%%%%%%%%
\begin{frame}[fragile]\frametitle{Chain-of-Thought}

\textbf{Prompt:} \emph{Explain how to make a peanut butter and jelly sandwich. Let's think step by step.}

\medskip
\textbf{What changes:} Adding \textbf{``Let's think step by step''} triggers the AI to slow down and reason through the problem sequentially instead of jumping straight to an answer.

\medskip
\textbf{Lesson:} This tiny addition is one of the most powerful prompt engineering tricks ever discovered. It works especially well for problems that require \textbf{reasoning, planning, or multi-step logic}. The AI makes fewer mistakes when it ``shows its work''.

\note{Chain-of-Thought (CoT) was formalised by Wei et al. (2022). Adding just ``Let's think step by step'' to a prompt can lift accuracy on arithmetic and logic tasks by 20-40\% on benchmark datasets.}

\end{frame}

%%%%%%%%%%%%%%%%%%%%%%%%%%%%%%%%%%%%%%%%%%%%%%%%%%%%%%%%%%%
\begin{frame}[fragile]\frametitle{Self-Consistency}

\textbf{Prompt:} \emph{Here are recipes of multiple sandwiches: Sandwich 1: [recipe 1]. Sandwich 2: [recipe 2]. \ldots\ Explain how to make a peanut butter jelly sandwich, then pick the most consistent approach across all variations.}

\medskip
\textbf{What changes:} The AI considers multiple paths to the same answer, then picks the most reliable one, like asking several people the same question and going with the majority view.

\medskip
\textbf{Lesson:} \textbf{Self-Consistency Prompting} is useful when accuracy matters more than speed. Instead of trusting one answer, you generate several reasoning paths and pick the one most answers agree on.

\end{frame}

%%%%%%%%%%%%%%%%%%%%%%%%%%%%%%%%%%%%%%%%%%%%%%%%%%%%%%%%%%%
\begin{frame}[fragile]\frametitle{Reflection and Iteration, The Big Picture}

\begin{itemize}
\item Every slide in this section illustrated one prompt technique applied to the \textbf{same task} (making a sandwich).
\item \textbf{Same AI model, same task, very different results}, only the prompt changed.
\end{itemize}

\medskip
\begin{center}
\renewcommand{\arraystretch}{1.3}
\begin{tabular}{|l|l|}
\hline
\textbf{Technique Used} & \textbf{What It Controls} \\
\hline
Role / Persona & Tone and expertise level \\
Constraints & Scope and safety \\
Few-Shot Examples & Output format and style \\
Context & Relevance and specificity \\
Feedback / Iteration & Quality improvement \\
Chain-of-Thought & Accuracy on complex tasks \\
Self-Consistency & Reliability and correctness \\
\hline
\end{tabular}
\end{center}

\medskip
\textbf{Key Takeaway:} Prompt Engineering is not magic, it is structured communication.

\end{frame}
